\begin{textsample}{NEG Dim 1 – rr – Score 3.00 – rr/rr\_em\_44.txt}  \label{ex:f1_neg_028}
Os Itinerários Formativos são organizados a partir de quatro eixos estruturantes : Investigação Científica, Processos Criativos, Mediação e Intervenção Sociocultural e Empreendedorismo.
Tais eixos visam integrar e integralizar os diferentes arranjos de Itinerários Formativos, bem como criar oportunidades para que os estudantes vivenciem experiências educativas profundamente associadas à realidade contemporânea, que promovam a sua formação pessoal, profissional e cidadã ( BRASIL, 2018d ).
Como os quatro eixos estruturantes são complementares, é importante que os Itinerários Formativos incorporem e integrem todos eles, a fim de garantir que os estudantes experimentem diferentes situações de aprendizagem e desenvolvam um conjunto diversificado de habilidades relevantes para sua formação integral.
Para tanto, é fundamental envolver os alunos em situações de aprendizagem que os permitam produzir conhecimentos, criar, intervir na realidade e empreender projetos presentes e futuros.
Para possibilitar melhor compreensão, apresentaremos os quatro eixos estruturantes, conforme os Referenciais Curriculares para a Elaboração de Itinerários Formativos – Portaria nº 1.432 ( BRASIL, 2018h ).

% matched lemmas: 
\end{textsample}
